% /root/monodromy, 19 November 2026.
% Elementary all-dimension characteristic-two computation; root independently audited.
\subsection{Hypersurfaces in characteristic two: an elementary exact group calculation}
\label{subsec:pn-char2-exact}
\status{Proved; independently audited.}

Fix an algebraically closed field $k$ of characteristic two,
$n\geq2$, and $d\geq n+2$.  Put $m=d-1$ and
$h=\binom{m}{n}$.  Work on the full Hasse--Witt invertibility
open $W_{n,d}$ as in the preceding section.  Its Lang fibre
at a general-position hyperplane arrangement is
$V=\bigwedge^n W$, where
$W=\mathbf F_2^d/\langle e_0+\cdots+e_m\rangle$.
Use $e_1,\ldots,e_m$ as a basis of $W$, so
$e_0=e_1+\cdots+e_m$.
Let $b_A=\bigwedge_{i\in A}e_i$, indexed by the $n$-subsets
$A\subset\{1,\ldots,m\}$, be the exterior basis of $V$.
Signs disappear in this characteristic.

\begin{lemma}\label{lem:char2-independent-quadric-swaps}
The geometric monodromy contains, for each
$(n-1)$-subset $I\subset\{2,\ldots,m\}$, the transvection
\begin{equation}\label{eq:char2-quadric-transvection}
 R_I=1+w_I\phi_I,\qquad
 w_I=\sum_{\substack{a\in\{2,\ldots,m\}\\a\notin I}}b_{I\cup\{a\}},
 \quad
 \phi_I=b_{\{1\}\cup I}^*.
\end{equation}
It fixes every exterior basis vector except
$b_{\{1\}\cup I}$, which it sends to
$\sum_{a\notin I}b_{I\cup\{a\}}$.
The swaps for different $I$ are independent monodromy elements.
\end{lemma}

\begin{proof}
Replace $H_0H_1$ by a varying quadric and keep the remaining
hyperplanes fixed.  The same alternating resolution as in
Lemma~\ref{lem:quadric-reflections} has constant lower terms
and a pair of points for every line
$L_I=\bigcap_{i\in I}H_i$.  On an open containing the original
arrangement these pairs are distinct, all positive-dimensional
quadric strata are connected and have no higher coherent
cohomology, and the entire top coherent Frobenius is bijective.

We check independence of the point-pair covers in characteristic
two, working on the affine space of quadric equations.  This
cover is the pullback of the projective parameter cover, so
its full image is already a sufficient geometric lower bound.
On a chosen basis for a line $L$, the restricted quadratic
form is $a_LX^2+b_LXY+c_LY^2$.  Its ordered roots give the
Artin--Schreier extension
\[
 z^2+z=a_Lc_L/b_L^2
\]
of the rational function field of the space of quadrics.
The restriction map to any line is surjective, so
$a_L,b_L,c_L$ are independent linear functions.
Moreover $b_L$ is evaluation of the polar alternating form
on the Pluecker vector of $L$.  For distinct lines these
linear functions are nonassociate, by injectivity of the
Pluecker embedding.

At the prime divisor $b_{L_I}=0$, the displayed function for
$I$ has pole order two and leading residue $a_{L_I}c_{L_I}$.
This residue is not a square in the divisor's rational
function field, since $a_{L_I},c_{L_I}$ remain independent
variables there.  The corresponding functions for all
$J\ne I$ have no pole at this divisor.
No Artin--Schreier coboundary $q^2+q$ can have that leading
term: a pole of order two would require a pole of order one
for $q$ and would have square leading residue.  Thus no
nonempty sum of these functions is an Artin--Schreier
coboundary.  Their classes are independent, and the compositum
has Galois group $(\mathbf Z/2)^{\binom{d-2}{n-1}}$.
Passing to the nonempty ordinary and transverse open preserves
connectedness of this cover and hence all the independent swaps.

At the arrangement, swapping the two points on $L_I$ exchanges
the cochains $e_0\wedge e_I$ and $e_1\wedge e_I$, and fixes
all other top cochains not involving both $0$ and $1$.
These latter cochains span the arrangement cohomology; those
involving both are eliminated by the coboundary relation
$u\wedge\bigwedge^{n-1}P=0$.
In the chosen exterior basis the only changed column is
$b_{\{1\}\cup I}$, and its difference is
$(e_0+e_1)\wedge e_I=w_I$.
This proves \eqref{eq:char2-quadric-transvection}.
The vector $w_I$ is nonzero because
$d-n-1\geq1$, and its $\phi_I$ value is zero.
Thus the operator is a transvection.
\end{proof}

\begin{lemma}\label{lem:char2-coordinate-generation}
Let $G_0$ be generated by the arrangement permutation group
and the swaps of Lemma~\ref{lem:char2-independent-quadric-swaps}.
Writing $\ell=\sum_A b_A^*$, one has
\[
 G_0=
 \begin{cases}
 \mathrm{GL}_{h}(\mathbf F_2),&d-n-1\ \text{odd},\\
 \operatorname{Stab}_{\mathrm{GL}_h(\mathbf F_2)}(\ell),
       &d-n-1\ \text{even}.
 \end{cases}
\]
In the second case the stabilizer is
$\mathbf F_2^{h-1}\rtimes\mathrm{GL}_{h-1}(\mathbf F_2)$.
\end{lemma}

\begin{proof}
Permuting the labels $1,\ldots,m$ and fixing label $0$ supplies
the following operators for every $(n-1)$-subset $I$ and
every $j\notin I$:
\[
 T^I_j(b_{I\cup\{j\}})=S_I,\qquad
 S_I=\sum_{a\notin I}b_{I\cup\{a\}},
\]
with every other exterior basis vector fixed.  This operator
exchanges the two vectors $b_{I\cup\{j\}}$ and $S_I$ and
fixes each of the other vectors in this sum.
For $j,t\notin I$, $j\ne t$, the conjugate
$T^I_jT^I_tT^I_j$ is therefore the pure coordinate
transposition exchanging $b_{I\cup\{j\}}$ and
$b_{I\cup\{t\}}$ and fixing all other exterior basis vectors.
The graph of $n$-subsets joined when they share $n-1$ elements
is connected.  Its edge transpositions generate the full
coordinate symmetric group $S_h$, so $G_0$ contains that group.

Put $l=d-n-1$.  Each $T^I_j$ is a transvection whose image
vector has exactly $l$ nonzero exterior coordinates, all
different from its input coordinate.  Conjugating by $S_h$
gives every transvection
$1+v b_a^*$ with $v$ supported on any $l$ coordinates other
than $a$.  If $l=1$, these are all elementary transvections
and generate $\mathrm{GL}_h(\mathbf F_2)$.
If $l\geq2$, choose two such supports with $l-1$ common
coordinates and differing in any desired pair $b,c$, both
different from $a$.  This is possible because
$h=\binom{d-1}{n}\geq d-1\geq d-n+1=l+2$
for $n\geq2,d\geq n+2$.
The product, with the same input functional, is
\[
 1+(b_b+b_c)b_a^*.
 \tag{4}
\]
If $l$ is odd, adding pairs from an $l$-element support reduces
it to a singleton.  Thus once again all elementary
transvections occur and $G_0=\mathrm{GL}_h(\mathbf F_2)$.

Suppose $l$ is even.  Every generator fixes $\ell$:
the coordinate transvections change a column by an even-weight
vector, and an arrangement permutation fixes $\ell$ as well.
For the latter assertion, permutations of $1,\ldots,m$ act
by coordinate permutations, while $(01)$ is the product of
the commuting swaps in
Lemma~\ref{lem:char2-independent-quadric-swaps}.
These permutations generate the full arrangement $S_d$.
Consequently $G_0\subseteq\operatorname{Stab}(\ell)$.

Relabel the exterior basis as $b_1,\ldots,b_h$ and put
$f_i=b_i+b_h$ for $1\leq i<h$.
The $f_i$ form a basis of $\ker\ell$, with complement
$\langle b_h\rangle$.  Operators (4), with image $f_i$ and
input functional $b_j^*$ for $i\ne j<h$, give the elementary
matrices on $\ker\ell$ and fix $b_h$.  They generate its
full $\mathrm{GL}_{h-1}(\mathbf F_2)$.
The coordinate transposition $(i\,h)$ is
$1+f_i(b_i^*+b_h^*)$.  Multiplying it by the commuting
operators $1+f_i b_j^*$ for all $j\ne i,h$ gives
$1+f_i\ell$.  These generate all the translations of the
complement by $\ker\ell$.  Hence the full stabilizer is
contained in $G_0$, proving equality.
\end{proof}

\begin{lemma}\label{lem:pn-char2-canonical-covector}
If $d\equiv n+1\pmod2$ and $d\geq n+1$, the universal ordinary
Hasse--Witt local system has a nowhere-zero invariant
$\mathbf F_2$-valued covector, defined over $\mathbf F_2$.
\end{lemma}

\begin{proof}
Write a degree-$d$ equation as $F=\sum_{|\alpha|=d}c_\alpha x^\alpha$.
Use the standard top-cohomology basis indexed by positive
$(n+1)$-tuples $u$ with $|u|=d$.  The Hasse--Witt matrix is
$H_{uv}=c_{2v-u}$, with an out-of-range coefficient understood
as zero.  Put
\[
 w_u=\begin{cases}c_u,&\text{every coordinate of }u\text{ is odd},\\
                   0,&\text{otherwise}.\end{cases}
\]
Then
\[
 H^{\mathsf T}w=w^{[2]}.
 \tag{5}
\]
Indeed, in the $v$ coordinate the sum is
$\sum_{u\ {\rm odd}}c_{2v-u}c_u$.  Whenever the coefficient
is nonzero, $2v-u$ is also a positive all-odd tuple.
The involution $u\leftrightarrow2v-u$ cancels all paired
terms.  Its only possible fixed term is $u=v$, giving
$c_v^2$ when $v$ is all odd and zero otherwise.

There exists a positive all-odd tuple of sum $d$ under the
assumed parity and bound.  If every $w_u$ vanished, the row
of $H$ indexed by any such tuple would vanish: every
possibly nonzero coefficient in that row is itself indexed
by an all-odd tuple.  Thus $w\ne0$ wherever $H$ is invertible.
The map $x\mapsto w^{\mathsf T}x$ is Frobenius equivariant
by (5), and restricts on Lang fixed vectors to an
$\mathbf F_2$-valued covector.  It is nonzero because those
fixed vectors span the coherent fibre.
This construction is global on projective parameter space:
the coherent bundle is $\mathcal O(-1)$ tensored with the
fixed top-cohomology vector space, and the entries $w_u$
are degree-one functions of the equation.  They define an
honest morphism of this bundle to $\mathcal O$, not merely
a covector up to scalar.  Everything is defined over
$\mathbf F_2$.
\end{proof}

\begin{theorem}\label{thm:pn-char2-all-groups}
Let $k$ be any perfect field of characteristic two and
$n\geq2$.  Put $h=\binom{d-1}{n}$.  The geometric and
arithmetic monodromy groups on the ordinary smooth
degree-$d$ hypersurface locus of $\mathbf P^n_k$ are:
\[
 \begin{array}{c|c}
 \text{degree range}&\text{both monodromy images}\\ \hline
 d\leq n&1\text{ on the zero vector space}\\
 d=n+1&\mathrm{GL}_1(\mathbf F_2)=1\\
 d\geq n+2,\ d\equiv n\pmod2&\mathrm{GL}_h(\mathbf F_2)\\
 d\geq n+3,\ d\equiv n+1\pmod2&
    \operatorname{Stab}(\lambda)
       \simeq\mathbf F_2^{h-1}\rtimes
                         \mathrm{GL}_{h-1}(\mathbf F_2),
 \end{array}
\]
where $\lambda$ is the nonzero covector in
Lemma~\ref{lem:pn-char2-canonical-covector}.
In the last case $\ker\lambda$ is a proper nonzero invariant
hyperplane, so the representation is reducible.
\end{theorem}

\begin{proof}
The two initial ranks are zero and one as stated.
In all higher degrees
Lemma~\ref{lem:char2-coordinate-generation} gives the required
geometric lower bound.  In the parity with a canonical
covector, Lemma~\ref{lem:pn-char2-canonical-covector} gives
the matching upper bound, even for the arithmetic image.
For precision in comparing covectors: the subgroup $G_0$
has exactly the one nonzero invariant covector $\ell$.
The canonical $\lambda$ is nonzero and $G_0$-invariant,
so its value on the arrangement fibre is $\ell$.
Thus the two stabilizers are the same.
In the other parity the lower bound is the entire general
linear group.  The result over arbitrary perfect $k$
follows by geometric connectedness of the corresponding
finite cover under algebraically closed extension, exactly
as in the odd-characteristic theorem.
\end{proof}

\begin{corollary}\label{cor:pn-all-p-sharp-residue-threshold}
For every finite field $k$, every $n\geq2$, and every
$d\geq n+2$, each nonempty open in the smooth degree-$d$
hypersurface parameter space of $\mathbf P^n_k$ contains
hypersurfaces, over suitable finite extensions of $k$,
with every point-count residue modulo the characteristic.
The threshold $n+2$ is exact: for $d\leq n$ the only
residue is $1$, and for $d=n+1$ the ordinary open misses $1$.
\end{corollary}

\begin{proof}
Odd characteristic was settled in
Corollary~\ref{cor:pn-odd-all-residues-sharp}.  In characteristic
two, the full general linear group in rank at least two
has both trace values.  In the covector-stabilizer case,
the trace is $1+\operatorname{Tr}(B)$ with
$B\in\mathrm{GL}_{h-1}(\mathbf F_2)$; here $h-1\geq2$,
so again both values occur.  The same finite Chebotarev
and coherent trace argument gives both point-count
residues in every open.  The zero- and one-rank failures
are as in the preceding corollary.
\end{proof}

\paragraph{Finite checks.}
The retained program
\path{code/projective_quadric_generators.py} constructs
these exact exterior matrices over $\mathbf F_2$.
It checks that every swap has rank-one difference and
order two, that their product is the hyperplane
transposition, and that the common fixed covector space
has dimension one exactly in the indicated parity.
For the tested ranks at most fifteen it additionally
enumerates all nonzero vector orbits, obtaining either
$2^h-1$ or $2^{h-1}-1$ and $2^{h-1}$ as predicted.
These finite checks support the formulas; the uniform
group-generation proof above is independent of sampling.
